% Research candidate only. Not Rethlas verified and not authorized for live integration.

\begin{proof}[Proof of \cref{prop:p062-polydisc-trace-separation-candidate}]
Fix \(a\in D^-_{\mathbf R}(K)\) and \(\boldsymbol\rho\) as in the
statement, and put
\[
 \Omega_{\boldsymbol\rho}
 :=\prod_{i=1}^n(-\log\rho_i,\infty).
\]
The vanishing of the distributional Lagerberg Hessian in
\eqref{eq:p062-polydisc-radial-trace-candidate} says, in the standard
coordinates on \(\Omega_{\boldsymbol\rho}\), that every distribution
\(\partial^2\tau_{a,\boldsymbol\rho}/\partial u_i\partial u_j\)
vanishes. Hence
\begin{equation}\label{eq:p062-polydisc-affine-trace-candidate}
 \tau_{a,\boldsymbol\rho}(u)=c+\sum_{i=1}^n\lambda_i u_i
\end{equation}
for some \(c,\lambda_1,\ldots,\lambda_n\in\mathbb R\). For completeness,
this standard distributional assertion can be seen by convolving on
relatively compact Euclidean balls: every convolution has identically
zero ordinary Hessian and is therefore affine; compatibility on
overlapping balls, followed by convergence of the convolutions to the
continuous function \(\tau_{a,\boldsymbol\rho}\), gives
\eqref{eq:p062-polydisc-affine-trace-candidate}.

Fix \(i\) and fix all \(u_j\), \(j\neq i\). As \(u_i\to\infty\), the
Gauss points in \eqref{eq:p062-polydisc-radial-trace-candidate} converge
to the Gauss point on the coordinate section \(T_i=a_i\) with the other
polyradii unchanged. Indeed, this follows directly by taking the limit
in the Gauss seminorm
\[
 \left|\sum_{\nu\in\mathbf N^n}c_\nu(T-a)^\nu\right|_{
                    \eta_{a,(\mathrm e^{-u_1},\ldots,\mathrm e^{-u_n})}}
 =\max_\nu |c_\nu|\mathrm e^{-\langle\nu,u\rangle}.
\]
More formally, the convergence follows first on polynomials and then on
the affinoid algebra of a fixed closed polydisc by uniform approximation.
The limiting point still belongs to \(D^-_{\mathbf R}\). Continuity of
\(h\) therefore shows that the left-hand side of
\eqref{eq:p062-polydisc-affine-trace-candidate} has a finite limit as
\(u_i\to\infty\). Thus \(\lambda_i=0\). Doing this for every \(i\)
shows that the trace is constant. Letting all \(u_i\to\infty\) gives
\begin{equation}\label{eq:p062-polydisc-radial-constant-candidate}
                   h(\eta_{a,\mathbf r})=h(a)
 \qquad (0<r_i<\rho_i).
\end{equation}

Now let \(a,b\in D^-_{\mathbf R}(K)\). The ultrametric inequality gives
\(|a_i-b_i|<R_i\). Choose \(r_i,\rho_i\) such that
\[
             |a_i-b_i|\leq r_i<\rho_i<R_i.
\]
Translation invariance of the Gauss norm gives
\(\eta_{a,\mathbf r}=\eta_{b,\mathbf r}\). Applying
\eqref{eq:p062-polydisc-radial-constant-candidate} at both centres gives
\(h(a)=h(b)\). Thus \(h\) is constant on the \(K\)-rational points.
Rigid points are dense in every strictly affinoid space, and, since \(K\)
is algebraically closed, the rigid points of a strict polydisc are exactly
its \(K\)-rational points. Applying this on an exhaustion by closed strict
subpolydiscs and using continuity proves that \(h\) is constant on all of
\(D^-_{\mathbf R}\). This last step includes all non-Abhyankar points; no
claim that such points belong to a moment skeleton is being made.

For \(D^+_{\mathbf R}\), the same argument gives
\eqref{eq:p062-polydisc-radial-constant-candidate} for \(r_i<R_i\).
Allowing some \(r_i=R_i\) follows by increasing the corresponding radii
to \(R_i\) and using continuity of both the Gauss section and \(h\). For
any \(a,b\in D^+_{\mathbf R}(K)\), choose
\(0<r_i\leq R_i\) with \(|a_i-b_i|\leq r_i\), taking \(r_i=R_i\) when
necessary. Then again
\(\eta_{a,\mathbf r}=\eta_{b,\mathbf r}\), so \(h(a)=h(b)\). Density of
rigid points and continuity finish the closed-polydisc case.
\end{proof}

\begin{proof}[Proof of \cref{thm:p062-polydisc-zero-current-candidate}]
Fix \(a\in D^-_{\mathbf R}(K)\) and an arbitrary
\(\boldsymbol\rho\) with \(0<\rho_i<R_i\). The ultrametric inequality
shows that the polydisc with centre \(a\) and polyradius
\(\boldsymbol\rho\) is contained in \(D^-_{\mathbf R}\). Set
\[
 A_{a,\boldsymbol\rho}
 :=\bigl\{x:0<|T_i(x)-a_i|<\rho_i\text{ for every }i\bigr\},
\]
\[
 F_a(x):=(-\log|T_i(x)-a_i|)_{i=1}^n.
\]
Translation by \(a\) identifies \(A_{a,\boldsymbol\rho}\) with the open
torus domain
\(\operatorname{trop}^{-1}(\Omega_{\boldsymbol\rho})\subset
(\mathbf G_{\mathrm m,K}^n)^{\an}\). In particular, \(F_a\) is a smooth
tropicalization, and its canonical skeleton is the one-sheeted Gauss
section
\[
 \iota_a(u)=\eta_{a,(\mathrm e^{-u_1},\ldots,\mathrm e^{-u_n})}.
\]

Let \(\beta\) be an arbitrary smooth compactly supported Lagerberg form
of bidegree \((n-1,n-1)\) on
\(\Omega_{\boldsymbol\rho}\). The tropicalization map of a split torus is
proper. Consequently \(F_a^*\beta\) is a compactly supported smooth test
form on \(A_{a,\boldsymbol\rho}\). We use it after restricting the zero
current to this open analytic domain, equivalently after extending the
test form by zero to \(D^-_{\mathbf R}\). Restriction of
\eqref{eq:p062-polydisc-zero-current-candidate} and the definition of the
differentials on currents give
\begin{align}
 0
 &=\bigl\langle \mathrm d'\mathrm d''[h],F_a^*\beta\bigr\rangle \\
 &=\int_{A_{a,\boldsymbol\rho}}
          h\,F_a^*(\mathrm d'\mathrm d''\beta).
             \label{eq:p062-current-pulltest-candidate}
\end{align}

Write the top-degree Lagerberg form as
\[
 \mathrm d'\mathrm d''\beta
 =q(u)\,\mathrm d'u_1\wedge\mathrm d''u_1\wedge\cdots
          \wedge\mathrm d'u_n\wedge\mathrm d''u_n
\]
for a smooth compactly supported function \(q\), with the fixed global
sign dictated by the chosen ordering convention. Put
\(\Gamma=-\log|K^\times|\), and choose a compact
\((\mathbf Z,\Gamma)\)-polytope \(\Delta\Subset
\Omega_{\boldsymbol\rho}\) containing \(\operatorname{supp}(q)\). Put
\[
 \alpha:=\mathrm d'\mathrm d''\beta,\qquad
 \mu_\alpha:=(F_a^*\alpha)_{\mathrm{Rad}}.
\]
The support theorem for top-degree forms on a torus
(Burgos Gil--Gubler--Jell--K\"unnemann, Proposition~2.3.6) gives
\[
 \operatorname{supp}(\mu_\alpha)
 \subseteq\iota_a(\operatorname{supp}(\alpha)).
\]
The calibrated integration formula on the compact torus domain
\(F_a^{-1}(\Delta)\), whose tropical lattice multiplicity is one, gives
\[
            (F_a)_*\mu_\alpha=\alpha_{\mathrm{Rad}}.
\]
This is the split-torus case of Gubler--Jell--Rabinoff,
Proposition~11.1 and Example~11.7. Indeed, the formula first gives this
equality against compactly supported
smooth multipliers; their uniform density on the compact support and
uniqueness of signed Radon measures give the displayed measure identity.
Since \(F_a\) restricts on the canonical skeleton to a homeomorphism with
inverse \(\iota_a\), the support and pushforward identities imply
\[
             \mu_\alpha=(\iota_a)_*\alpha_{\mathrm{Rad}}.
\]
Consequently the integration formula holds for the arbitrary continuous
multiplier \(h\), with no approximation of \(h\) and no pullback of its
current. Thus
\eqref{eq:p062-current-pulltest-candidate} is exactly
\[
 0=\int_{\Omega_{\boldsymbol\rho}}
       h(\iota_a(u))\,\mathrm d'\mathrm d''\beta.
\]
It follows that the trace in
\eqref{eq:p062-polydisc-radial-trace-candidate} has zero distributional
Lagerberg Hessian. Since this holds for every centre and every
\(\boldsymbol\rho\) allowed in the statement,
\cref{prop:p062-polydisc-trace-separation-candidate} shows that \(h\) is
constant. A constant is weakly smooth and has zero weakly smooth Hessian;
Gubler--Jell--Rabinoff's characterization then makes it harmonic. It and
its negative pull back to harmonic functions on every analytic curve after
every extension of scalars, so it is classically pluriharmonic.

For the closed polydisc, every centred punctured polydisc with
\(\rho_i<R_i\) is an ordinary open analytic subspace of
\((\mathbf A_K^n)^{\an}\), and the ultrametric inequality places it inside
\(D^+_{\mathbf R}\). Hence it lies in the topological interior of this
analytic domain. Berkovich's relative-interior theorem for analytic
domains (\emph{\'Etale cohomology for non-Archimedean analytic spaces},
Proposition~1.5.5(i)) identifies that topological interior with the
relative analytic interior; since analytic affine space is boundaryless,
the composition formula for interiors places the punctured polydisc in
the absolute analytic interior of \(D^+_{\mathbf R}\). The identical test-form argument
therefore proves the trace hypothesis there, and the closed case of
\cref{prop:p062-polydisc-trace-separation-candidate} applies.
\end{proof}
